\documentclass{article}[12pt]
\usepackage{graphicx}
\usepackage{amsmath,amssymb}
\usepackage{tikz}
\usepackage{xepersian}
\settextfont[Scale=1]{IRXLotus}
\setlatintextfont[Scale=0.8]{}

\DeclareRobustCommand{\bbone}{\text{\usefont{U}{bbold}{m}{n}1}}

\DeclareMathOperator{\EX}{\mathbb{E}}% expected value
\let\P\relax
\DeclareMathOperator{\P}{\mathbb{P}}

\title{  \includegraphics[scale=0.35]{../logo.png} \\
    دانشکده مهندسی کامپیوتر
    \\
    دانشگاه صنعتی شریف
}
\author{استاد درس: دکتر حمیدرضا ربیعی}
\date{پاییز ۱۴۰۱}



\def \Subject {
تمرین اول
}
\def \Course {
درس فرآیند تصادفی
}
\def \Author {
نام و نام خانوادگی:
امیر پورمند}
\def \Email {\lr{pourmand1376@gmail.com}}
\def \StudentNumber {99210259}


\begin{document}

 \maketitle
 
\begin{center}
\vspace{.4cm}
{\bf {\huge \Subject}}\\
{\bf \Large \Course}
\vspace{.8cm}

{\bf \Author}

\vspace{0.3cm}

{\bf شماره دانشجویی: \StudentNumber}

\vspace{0.3cm}

آدرس ایمیل
:
{\bf \Email}
\end{center}


\clearpage
\section{سوال ۱}
\subsection{الف}

زمان ارسال ایمیل برای علی تابع متغیر تصادفی 
$X$
خواهد بود که از توزیع پواسون با پارامتر لاندا پیروی میکند. 

\begin{equation}\begin{split}
	X \sim Pois(\lambda_A) \\
	P(X == 3) = \frac{\lambda^k e^{-\lambda}}{k!} = 
	\frac{\lambda_A^3 e^{-\lambda_A}}{3!}	 
	\end{split}
\end{equation}

\subsection{ب}
\subsubsection{الف}

\begin{equation}
\begin{split}
		f_{Y_{2} |Y_{1}}( y_{2} |y_{1}) =f_{X_{2}}( y_{2} -y_{1}) =\lambda e^{-\lambda ( y_{2} -y_{1})}\\
		E[ Y_{2} |Y_{1}] =\int _{0}^{\infty } Y_{2} f_{Y_{2} |Y_{1}}( y_{2} |y_{1}) dy_{2} =\int _{0}^{\infty } Y_{2} \lambda e^{-\lambda ( y_{2} -y_{1})} dy_{2}\\
		=\frac{e^{\lambda y_{1}}}{\lambda }
\end{split}
\end{equation}


\subsubsection{ب}

میدانیم تابع چگانی 
$Y_1$
بصورت زیر است:
\begin{equation}
	Y_{1} =\begin{cases}
		\lambda e^{-\lambda } & Y_{1} \geqslant 0\\
		0 & Y_{1} < 0
	\end{cases}
\end{equation}
سپس داریم
\begin{equation}
	\begin{split}
		 \begin{array}{c}
			X\ =\ Y_{1}^{2} \ \Longrightarrow \ Y_{1} \ =\ \sqrt{x}\\
			f_{X}( x) \ =\ f_{Y_{1}}\left(\sqrt{x}\right)\frac{1}{2\sqrt{x}} \ =\lambda e^{-\lambda \sqrt{x}}\frac{1}{2\sqrt{x}} \ \ 
		\end{array}
	\end{split}
\end{equation}
\subsubsection{ج}	

\begin{equation}
	\begin{split}
		  \begin{array}{c}
		 	f_{Y_{1} ,Y_{2}}( y_{1} ,y_{2}) =f_{Y_{2} |Y_{1}}( y_{2} |y_{1}) f_{Y_{1}}( y_{1})\\
		 	=f_{X_{1}}( y_{2} -y_{1}) f_{Y_{1}}( y_{1}) =\lambda e^{-\lambda ( y_{2} -y_{1})} \lambda e^{-\lambda y_{1}} =\lambda ^{2} e^{-\lambda y_{2}}
		 \end{array}
	\end{split}
\end{equation}


\subsection{ج}
\subsubsection{الف}

\begin{equation}
	\begin{gathered}
				E[ Y_{1} |Y_{1} \leqslant 1] =?\\
				F_{Y_{1} |Y_{1} \leqslant 1}( y_{1}) =P( Y_{1} \leqslant y_{1} |Y_{1} \leqslant 1) =1-P( Y_{1} \geqslant y_{1} |Y_{1} \leqslant 1)\\
				=1-\frac{P( y_{1} \leqslant Y_{1} \leqslant 1)}{P( Y_{1} \leqslant 1)} =1-\frac{e^{-\lambda y_{1}} e^{-\lambda ( 1-y_{1})} \lambda ( 1-y_{1})}{e^{-\lambda } \lambda }\\
				=1-( 1-y_{1}) =y_{1}
	\end{gathered}
\end{equation}

\begin{equation}
	\begin{gathered}
		E[ Y_{1} |Y_{1} \leqslant 1] =\int _{0}^{1}( 1-y_{1}) dy_{1} =\left[ y_{1} -y_{1}^{2} /2\right]_{0}^{1} =\frac{1}{2}
	\end{gathered}
\end{equation}

\subsubsection{ب}
\begin{equation}
	E[ Y_{2} |Y_{1} \leqslant 1] =E[ 1+X_{2}] =1+\frac{1}{\lambda }
\end{equation}

\subsection{د}
\subsubsection{الف}
خب فرض کنیم علی 
$A$
نامه میفرستد در حالی که محمد 
$M$
نامه میفرستد. و نرخ ارسال ایمیل های علی را 
$\lambda_1=2\lambda_A$
و نرخ ارسال ایمیل های محمد را 
$\lambda_2=\lambda_B$
در نظر میگیریم و توجه داریم که با توجه به این که زمان علی و محمد متفاوت است در واقع 

 توزیع جمع نامه ها برابر خواهد بود با:


\begin{equation}
	\begin{aligned}
		p_{Z}(z) &=P(Z=z) \\
		&=\sum_{j=0}^{z} P(X=j \& Y=z-j) \quad \text { so } X+Y=z \\
		&=\sum_{j=0}^{z} P(X=j) P(Y=z-j)  \\
		&=\sum_{j=0}^{z} \frac{e^{-\lambda_{1}} \lambda_{1}^{j}}{j !} \frac{e^{-\lambda_{2}} \lambda_{2}^{z-j}}{(z-j) !} \\
		&=\sum_{j=0}^{z} \frac{1}{j !(z-j) !} e^{-\lambda_{1}} \lambda_{1}^{j} e^{-\lambda_{2}} \lambda_{2}^{z-j} \\
		&=\sum_{j=0}^{z} \frac{z !}{j !(z-j) !} \frac{e^{-\lambda_{1}} \lambda_{1}^{j} e^{-\lambda_{2}} \lambda_{2}^{z-j}}{z !} \\
		&=\sum_{j=0}^{z}\left(\begin{array}{l}
			z \\
			j
		\end{array}\right) \frac{e^{-\lambda_{1}} \lambda_{1}^{j} e^{-\lambda_{2}} \lambda_{2}^{z-j}}{z !} \\
		&=\frac{e^{-\lambda}}{z !} \sum_{j=0}^{z}\left(\begin{array}{l}
			z \\
			j
		\end{array}\right) \lambda_{1}^{j} \lambda_{2}^{z-j} \\
		&=\frac{e^{-\lambda}}{z !}\left(\lambda_{1}+\lambda_{2}\right)^{z}
	\end{aligned}
\end{equation}
حال توزیع نهایی برابر توزیع پواسون با پارامتر جمع این دو یعنی 
$\lambda_1+\lambda_2 = 2\lambda_A+\lambda_B$
خواهد بود. 
\subsubsection{ب}

\begin{equation}
	\begin{gathered}
			S_{0} =min( T_{0} ,1) \Longrightarrow \\
			F_{S_{0}}( s_{0}) =P( min( T_{0} ,1) < s_{0})\\
			=1-P( min( T_{0} ,1)  >s_{0}) =1-P( s_{0} < T_{0} ,\ s_{0} < 1)\\
			=1-P( s_{0} < T_{0}) P( s_{0} < 1) =1-II( s_{0} < 1)( 1-F_{T_{0}}( s_{0}))\\
			=1-\left( 1-\left( 1-e^{-\lambda s_{0}}\right)\right) =1-e^{-\lambda s_{0}}\\
			E[ S_{0}] =\int _{0}^{1} e^{-\lambda s_{0}} ds_{0} =\frac{-1}{\lambda }\left[ e^{-\lambda s_{0}}\right]_{0}^{1} =\frac{-1}{\lambda }\left( e^{-\lambda } -1\right) =\frac{1-e^{-\lambda }}{\lambda }\\
			E[ Q] =E[ S_{0}] +E[ T_{1}] =E[ min( T_{0} ,1)] +E[ exp( \lambda _{A})] =\frac{1}{\lambda_A } +\frac{1-e^{-\lambda_A }}{\lambda_A }
	\end{gathered}
\end{equation}

\subsubsection{ج}

\begin{equation}
	\begin{gathered}
			P( A=4|A+B=10) =P( A=4,B=6|A+B=10)\\
			=\frac{P( A=4) P( B=6)}{F_{A+B}( 10)} =\frac{e^{-2\lambda _{A}}( 2\lambda _{A})^{4}}{4!} \times \frac{e^{-\lambda _{B}} \lambda _{B}^{6}}{6!} \times \frac{10!}{e^{-2\lambda _{A} -\lambda _{B}}( 2\lambda _{A} +\lambda _{B})^{10}}\\
			=\frac{10*9*8*7}{4*3*2} \times \frac{2^{4} \lambda _{A}^{4} \lambda _{B}^{6}}{( 2\lambda _{A} +\lambda _{B})^{10}}\\
			=\frac{10*3*7*16*\lambda _{A}^{4} \lambda _{B}^{6}}{( 2\lambda _{A} +\lambda _{B})^{10}}
	\end{gathered}
\end{equation}

\subsection{ه}
\begin{equation}
	\begin{gathered}
				Chebyshev\ \\
				X\ \sim \ Pois( 4)\\
				E[ X] =4,\ Var( X) =4,\ \sigma ( X) =2\\
				P( |X-\mu |\geqslant k\sigma ) \leqslant \frac{1}{k^{2}}\\
				P( X\geqslant 5)\rightarrow P( X-4\geqslant 1)\\
				k\sigma =2k\leqslant 1\rightarrow k\leqslant \frac{1}{2}\rightarrow \frac{1}{k^{2}} \geqslant 4\ \rightarrow useless
	\end{gathered}
\end{equation}

\begin{equation}
	\begin{gathered}
			Markov\\
			X\sim Pois( 4)\\
			E[ X] =4\\
			P( X\geqslant a) \leqslant \frac{E[ X]}{a}\\
			P( X\geqslant 5) \leqslant \frac{4}{5} =0.8\ \rightarrow more\ useful
	\end{gathered}
\end{equation}
\subsection{و}
برای این مسئله باید ابتدا قضیه CLT را بیان کنم که میگوید اگر یک سری متغیر iid داشته باشیم توزیع میانگین این متغیرها همواره یک توزیع نرمال است. به طور دقیق تر برای این مثال وقتی 
$\lambda$
به اندازه کافی بزرگ باشد یعنی به اندازه 
$\lambda$
متغیرتصادفی داریم که هر کدام توزیع پواسون با پارامتر یک دارند و البته واریانس آنها نیز یک است. در اینجا جمع این متغیرهای پواسون همواره یک توزیع نرمال خواهد بود و برای تخمین میتوان از ان استفاده کرد. حال میتوان به جای حساب کردن توزیع پواسون در بازه صفر تا یک از توزیع نرمال با پارامتر لاندا به عنوان میانگین و البته واریانس لاندا برای این منظور استفاده کرد.
$N \sim N(\lambda,\lambda)$
\subsection{ز}

\begin{eqnarray*}
	\begin{gathered}
				\Lambda \sim Exp( 2)\\
				N|\Lambda \sim Pois( \lambda )\\
				P( N=n) =\int _{0}^{\infty } P( N=n|\Lambda =\lambda ) P( \Lambda =\lambda ) d\lambda \\
				=\int _{0}^{\infty }\frac{e^{-\lambda } \lambda ^{n}}{n!} \times 2e^{-2\lambda } d\lambda \\
				we\ know\ that\ \int _{0}^{\infty }\frac{e^{-x} x^{n}}{n!} dx=1\\
				=\frac{2}{3*3^{n}}\int _{0}^{\infty }\frac{e^{-3\lambda } 3^{n} \lambda ^{n}}{n!} d\lambda =\frac{2}{3^{n+1}}\\
				E\left[ N^{2}\right] =\int _{0}^{\infty } n^{2}\frac{2}{3^{n+1}} dn=\frac{4}{3ln^{3}( 3)}
	\end{gathered}
\end{eqnarray*}


\section{سوال ۲}
\subsection{الف}
\begin{equation}
	\begin{gathered}
			x( t) \ *\ h( t) =\int _{-\infty }^{+\infty } x( \tau ) h( t-\tau ) d\tau =\int _{-\infty }^{+\infty } e^{-a\tau } u( \tau ) e^{-b( t-\tau )} u( t-\tau ) d\tau \\
			=\int _{0}^{t} e^{-a\tau -bt+b\tau } d\tau =e^{-bt}\frac{1}{b-a}\left[ e^{( b-a) \tau }\right]_{0}^{t} =\frac{e^{-bt}}{b-a}\left[ e^{bt-at} -1\right] ,\ t >0
	\end{gathered}
\end{equation}
\subsection{ب}

\begin{equation}
	\begin{gathered}
			x( t) \ *\ h( t) =\int _{-\infty }^{+\infty } x( \tau ) h( t-\tau ) d\tau =\int _{-\infty }^{+\infty } e^{-a\tau } u( \tau ) u( t-\tau ) d\tau \\
			=\int _{0}^{t} e^{-a\tau } d\tau =\frac{-1}{a}\left[ e^{-a\tau }\right]_{0}^{t} =\frac{e^{-at} -1}{-a} ,\ t >0
	\end{gathered}
\end{equation}
\subsection{ج}
\begin{equation}
	\begin{gathered}
			x( t) \ *\ h( t) =\int _{-\infty }^{+\infty } x( \tau ) h( t-\tau ) d\tau =\int _{-\infty }^{+\infty } u( \tau ) u( t-\tau ) d\tau \\
			=\int _{0}^{t} d\tau =t,t >0
	\end{gathered}
\end{equation}


\section{سوال ۳}

\begin{equation}
	\begin{gathered}
		h( n) =\left(\frac{1}{5}\right)^{n} u[ n] ,\ \begin{cases}
			causal & n< 0\Longrightarrow h[ n] =0\\
			stable & \sum _{k=-\infty }^{+\infty } |h[ n] |=1+\frac{1}{5} +\frac{1}{25} +...\ =\ \frac{1}{1-\frac{1}{5}} =\frac{5}{4} < \infty 
		\end{cases}
	\end{gathered}
\end{equation}

\begin{equation}
	\begin{gathered}
		h( n) =0.8^{n} u( n+2) ,\begin{cases}
			noncausal & n=1,\ h[ n] \neq 0\\
			stable & 0.8^{-2} +0.8^{-1} +...=\frac{0.8^{-2}}{1-0.8} < \infty 
		\end{cases}
	\end{gathered}
\end{equation}

\begin{equation}
	\begin{gathered}
		h( n) =\left(\frac{1}{2}\right)^{n} u( -n) ,\begin{cases}
			noncausal & u< 0,h[ n]  >0\\
			not\ stable & 0.5^{0} +0.5^{-1} +...\ =\ \infty 
		\end{cases}
	\end{gathered}
\end{equation}

\begin{equation}
	\begin{gathered}
		h( n) =5^{n} u( 3-n) ,\begin{cases}
			noncausal & n=-1,h( n) \neq 0\\
			stable & 5^{3} +5^{2} +...\ =\frac{5^{3}}{1-\frac{1}{5}}
		\end{cases} < \infty 
	\end{gathered}
\end{equation}


\begin{equation}
	\begin{gathered}
		h( t) =e^{-4t} u( t-2) ,\begin{cases}
			casual & t< 0,h( t) =0\\
			stable & \int _{2}^{\infty } e^{-4t} dt=\frac{-1}{4}\left[ e^{-4t}\right]_{2}^{\infty } < \infty 
		\end{cases}
	\end{gathered}
\end{equation}


\begin{equation}
	\begin{gathered}
		h( t) =e^{-6t} u( 3-t) ,\begin{cases}
			noncausal & t=-1,h( t) \neq 0\\
			nonstable & \int _{-\infty }^{3} e^{-6t} dt=\infty 
		\end{cases}
	\end{gathered}
\end{equation}


\begin{equation}
	\begin{gathered}
		h( t) =e^{-2t} u( t+50) ,\begin{cases}
			noncausal & t=-1,h( t) \neq 0\\
			stable & \int _{-50}^{\infty } e^{-2t} dt=\frac{-1}{2}\left[ e^{-2t}\right]_{-50}^{\infty } < \infty 
		\end{cases}
	\end{gathered}
\end{equation}

\section{سوال ۴}
\begin{equation}
	\begin{gathered}
			y( t) =x( t-2) +x( 2-t)\\
			h( t) =\delta ( t-2) +\delta ( 2-t)\\
			stable:\ \int _{-\infty }^{+\infty } \delta ( t-2) +\delta ( 2-t) dt=2\\
			casual:\ t< 0,h( t) =0\\
			not\ memoryless:\ h( 2) =2\neq 0\\
			Linear:\ \\
			y( t) =x_{1}( t-2) +x_{1}( 2-t) +x_{1}( t-2) +x_{1}( 2-t)\\
			y_{1}( t) =x_{1}( t-2) +x_{1}( 2-t)\\
			y_{2}( t) =x_{1}( t-2) +x_{1}( 2-t)\\
			\\
			y( t) =ax( t-2) +ax( 2-t) =ax( t-2) +x( 2-t) =ay_{a}( t)\\
			we\ could\ also\ easily\ say\ that\ sum\ of\ two\ linear\ systems\ is\ linear!\\
			\\
			Time\ invariant:\ y( t-k) =x( t-k-2) +x( 2-t+k)\\
			S[ x( t-k)] =x( t-k-2) +x( 2-t+k) \ \rightarrow both\ equal
	\end{gathered}
\end{equation}\hrule
\begin{equation}
	\begin{gathered}
			y( t) =cos( 3t) x( t)\\
			h( t) =cost( 3t) \delta ( t)\\
			stable:\ \int cos( 3t) \delta ( t) dt=1\\
			memoryless:\ h( t) =0\ execpt\ t=0\\
			casual:h( t) =0,t< 0\\
			linear:\ cos( 3t) \ is\ constant( not\ dependant\ on\ x( t)) \ everytime\ and\ product\ of\ constant\ signal\ \\
			to\ a\ linear\ signal\ is\ definitely\ linear.\ \\
			\\
			Time\ variant:\ \\
			y( t-k) =cos( 3t-3k) x( t-k) \neq cos( 3t) x( t-k)
	\end{gathered}
\end{equation}
\hrule
\begin{equation}
	\begin{gathered}
			y( t) =\int _{-\infty }^{2t} x( t) dt\\
			h( t) \ =\ u( t)\\
			not\ stable:\ \int u( t) =\infty \\
			causal:\ t< 0\rightarrow h( t) =0\\
			not\ memoryless:\ t=1,\ h( t) =1\\
			linear:\ \\
			y( t) =\int _{-\infty }^{2t} x_{1}( t) +x_{2}( t) dt=\int _{-\infty }^{2t} x_{1}( t) dt+\int _{-\infty }^{2t} x_{2}( t) dt\\
			linear\ because\ integral\ is\ separable\\
			also\ integral\ can\ carry\ contant\ out\ \\
			\\
			Time\ variant:\ \\
			y( t) =\int _{-\infty }^{2t} x( t-k) dt\neq \int _{-\infty }^{2( t-k)} x( t-k) dt
	\end{gathered}
\end{equation}
\hrule
\begin{equation}
	\begin{gathered}
			y( t) =x( t/3)\\
			h( t) =\delta ( t/3)\\
			stable:\ \int \delta ( t/3) dt=1\\
			causal:\ t< 0,\ h( t) =0\\
			memoryless:h( t) =0\ except\ t=0\\
			linear\ since\ any\ linear\ signal\ with\ scaling\ remains\ linear\\
			time\ invariant\ since\ linear\ system\ x( t) \ is\ only\ scaled
	\end{gathered}
\end{equation}
\hrule
\begin{equation}
	\begin{gathered}
			y( t) =\begin{cases}
				x( t) +x( t-2) & t\geqslant 0\\
				o & ow
			\end{cases}\\
			h( t) =\ \begin{cases}
				\delta ( t) +\delta ( t-2) & t\geqslant 0\\
				0 & ow
			\end{cases}\\
			stable:\ \int \delta ( t) +\delta ( t-2) dt< \infty \\
			causal:\ t< 0,\ h( t) =0\\
			not\ memoryless:\ h( 2) \neq 0\\
			linear:\ \\
			\begin{cases}
				x_{1}( t) +x_{1}( t-2) +x_{2}( t) +x_{2}( t-2) & t\geqslant 0\\
				0 & ow
			\end{cases} =\\
			\begin{cases}
				x_{1}( t) +x_{1}( t-2) & t\geqslant 0\\
				o & ow
			\end{cases} +\begin{cases}
				x_{2}( t) +x_{2}( t-2) & t\geqslant 0\\
				o & ow
			\end{cases}\\
			it\ is\ the\ same\ if\ we\ scale\ it\ with\ some\ variable\ or\ give\ it\ the\ system\ of\ \\
			ax( t) .\ \\
			\\
			time\ variant:\\
			\begin{cases}
				x( t-k) +x( t-k-2) & t\geqslant k\\
				0 & ow
			\end{cases} \neq \begin{cases}
				x( t-k) +x( t-k-2) & t\geqslant 0\\
				0 & ow
			\end{cases}
	\end{gathered}
\end{equation}
\hrule
\begin{equation}
	\begin{gathered}
			y( t) =\begin{cases}
				x( t) +x( t-2) & x( t) \geqslant 0\\
				o & ow
			\end{cases}\\
			h( t) =\ \begin{cases}
				\delta ( t) +\delta ( t-2) & \\
				0 & 
			\end{cases}\\
			stable:\ \int \delta ( t) +\delta ( t-2) dt< \infty \\
			causal:\ t< 0,\ h( t) =0\\
			not\ memoryless:\ h( 2) \neq 0\\
			not\ linear:\\
			\begin{cases}
				x_{1}( t) +x_{1}( t-2) +x_{2}( t) +x_{2}( t-2) & x_{1}( t) +x_{2}( t) \geqslant 0\\
				o & ow
			\end{cases} \neq \\
			\begin{cases}
				x_{1}( t) +x_{1}( t-2) & x_{1}( t) \geqslant 0\\
				o & ow
			\end{cases} +\begin{cases}
				x_{2}( t) +x_{2}( t-2) & x_{2}( t) \geqslant 0\\
				o & ow
			\end{cases}\\
			\\
			time\ invariant:\\
			\begin{cases}
				x( t-k) +x( t-k-2) & x( t-k) \geqslant 0\\
				o & ow
			\end{cases} =\begin{cases}
				x( t-k) +x( t-k-2) & x( t-k) \geqslant 0\\
				o & ow
			\end{cases}
	\end{gathered}
\end{equation}

\section{سوال ۵}

\begin{equation}
	\begin{gathered}
			R_{x}( \omega ) =\int _{-\infty }^{+\infty } \delta ( w-w_{0}) +\delta ( w+w_{0}) e^{2j\pi \omega \tau } dw=e^{-2j\pi \omega _{0} \tau } +e^{2j\pi \omega _{0} \tau }\\
			R_{x}( \omega ) =\int _{-\infty }^{+\infty } e^{\frac{-\omega ^{2}}{2}} e^{2j\pi \omega \tau } dw=\sqrt{2\pi } \ e^{2j^{2} \pi ^{2} \tau ^{2}}\\
			R_{x}( w) =\int _{-\infty }^{+\infty } e^{-|\omega |} e^{2j\pi \omega \tau } dw\\
			=\int _{0}^{+\infty } e^{-\omega } e^{2j\pi \omega \tau } dw+\int _{-\infty }^{0} e^{\omega } e^{2j\pi \omega \tau } dw\\
			=\left[\frac{1}{2j\pi \tau -1} e^{w( 2j\pi \tau -1)}\right]_{0}^{\infty } +\left[\frac{1}{2j\pi \tau +1} e^{w( 2j\pi \tau +1)}\right]_{-\infty }^{0}\\
			=\frac{-1}{2j\pi \tau -1} +\frac{1}{2j\pi \tau +1}
	\end{gathered}
\end{equation}


 \end{document}